\section{Introduction}
LLM inference workloads differ sharply from training: decoding operates at batch
sizes of 1--8 with skinny matrices and tight latency budgets. These shapes are
often memory-bound, and naive quantization can underperform if dequantization
overhead is not amortized. This work explores packed INT4 GEMM in Triton and
provides a systems-oriented evaluation emphasizing utilization, bottleneck
regimes, and deployment-relevant shapes.

We argue that decode inference occupies a distinct optimization regime: launch
overhead and memory traffic dominate, so quantization benefits depend on
amortizing dequantization cost rather than peak compute.

Across representative decode shapes, Tiny-GEMM achieves up to 3.7$\times$
speedup over FP16 on wide FFN projections (M=1, K=4096, N=14336).

\paragraph{Contributions.}
\begin{itemize}
  \item A decode-focused INT4 GEMM kernel with static configuration selection
        and measurement-driven tuning.
  \item A regime analysis of INT4 decode GEMMs that identifies when quantization
        helps or hurts, supported by FP16 and dequantized FP16 baselines.
  \item A causal attribution methodology combining hardware counters and
        microbenchmarks to explain performance transitions.
\end{itemize}
